\documentclass[10pt,a4paper]{amsart}
\usepackage[margin=19mm]{geometry}
\usepackage{amsmath,amssymb,mathtools}
\usepackage[scr=rsfs,cal=euler]{mathalfa}
\usepackage{xfrac}
\usepackage[unicode,hidelinks,bookmarks=false]{hyperref}
\newcommand{\ie}{i.e.\ }
\newcommand{\cf}{cf.\ }
\newcommand{\resp}{respectively\ }
\newcommand{\Subsection}{\subsection}
\providecommand{\nobreakdash}{\nobreak\hbox{-}}
\setlength{\emergencystretch}{2em}
\multlinegap0pt
\usepackage{amssymb}
\usepackage[shortlabels]{enumitem}
\usepackage{tikz}
\newtheorem{lemma}{Lemma}
\def\Hfour{H_4}
\def\Hom{\mathop{\mathrm{Hom}}}
\def\Rset{\mathbb{R}}
\newcommand{\abs}[1]{\left\vert#1\right\vert}
\newcommand{\cupdot}{\mathbin{\mathaccent\cdot\cup}}
\def\nHom{\textsc{\#Hom}}
\def\nPolytopeVertices{\textsc{\#Vertices of a 0/1 Polytope}}
\title{Approximate counting of vertices of 0/1 polytopes: \\ a stronger hardness result}
\author{Heng Guo, Mark Jerrum}
\date{Source: arXiv:2608.22290v2 (CC BY 4.0)}
\begin{document}
\maketitle

\noindent\emph{Source and licence.} Approximate counting of vertices of 0/1 polytopes: \\ a stronger hardness result. Version arXiv:2608.22290v2 (CC BY 4.0), distributed by arXiv under the Creative Commons Attribution 4.0 International licence, http://creativecommons.org/licenses/by/4.0/, as recorded in the arXiv licence metadata for this exact version and shown on the abstract page https://arxiv.org/abs/2608.22290v2. Full licence notice: \texttt{licenses/arxiv-2608.22290-CC-BY-4.0.txt}.\par\smallskip

\noindent\textit{Editorial scope.} Counted unit: the article's lemma lem:reduction (source0.tex lines 170--172). The article's complete original proof of it is delivered with it (source0.tex lines 174--215, one \texttt{proof} environment of 42 lines). No mathematics is written, repaired or re-derived: the statement and the proof are the article's own lines, in the article's own notation, and no retained line is substituted or re-flowed.  No further source-local context is needed: every cross-reference of every retained line resolves inside the two delivered spans.  Nothing was omitted from any retained span, so the package carries no omission record.  No retained line contains a citation, so the wrapper carries no bibliography and no bibliography entry.  No retained line contains a figure environment, an image-inclusion command or drawing code, so no image is delivered and the package declares no asset dependency.  Editorial formatting only, in addition: an AMS \texttt{amsart} preamble that restates the article's own declarations and macros used by the retained lines, namely the environments \texttt{lemma} on separate counters, together with the standard packages those lines need.  The retained environments print in this document as Lemma 1 in order of appearance, whereas the article prints the same items with the numbering of its own full document.  The complete native source of the article as distributed in the arXiv e-print for this exact version is bundled unchanged as \texttt{sources/208322/source0.tex}.
\begin{lemma}\label{lem:reduction}
$\nHom(\Hfour)$ is AP-reducible to $\nPolytopeVertices$.
\end{lemma}

\begin{proof}
Let $G$ be an instance of $\nHom(\Hfour)$.  
We construct a polytope $Q(G)$ in a $2|V(G)|$-dimensional ambient space.  Let $y\in\mathbb{R}^{V(G)\times\{1,2\}}$ be a vector,
and for every $v\in V(G)$ write
\[
 s_v=y_{v,1}+y_{v,2}.
\]
Define the polytope $Q(G)\subseteq\Rset^{2\abs{V(G)}}$ as the set of vectors $y$ satisfying
\begin{align}
 0\leq y_{v,i}\leq1, &\quad\text{for all $v\in V(G)$ and $i\in\{1,2\}$}; \label{eq:box}\\
 s_u+s_v\geq2,&\quad\text{for all $\{u,v\}\in E(G)$}. \label{eq:edge}
\end{align}
Note that $(s_v)$ are auxiliary variables introduced for convenience in the proof, so \eqref{eq:edge} should be interpreted as a shorthand for 
$$
y_{u,1}+y_{u,2}+y_{v,1}+y_{v,2}\geq2.
$$
We aim to show that the vertices of $Q(G)$ are in bijection with homomorphisms from $G$ to~$\Hfour$.  This gives the required AP-reduction from $\nHom(\Hfour)$ to $\nPolytopeVertices$.
The key observation is that every vertex of $Q(G)$ belongs to $\{0,1\}^{2\abs{V(G)}}$.

Suppose, for a contradiction, that $y$ is a vertex of $Q(G)$ with a fractional coordinate.  We call $(y_{v,1},y_{v,2})$ the \textit{block} corresponding to~$v$.
We first claim that every fractional block has exactly one fractional coordinate.  Indeed, suppose both coordinates in the block corresponding to~$v$ are fractional.  For all sufficiently small $\delta>0$, we can add $\delta$ to one coordinate and subtract $\delta$ from the other.  This preserves $s_v$, and hence preserves all the inequalities in \eqref{eq:edge}.  Both directions of the perturbation satisfy the box constraints in \eqref{eq:box}.  This contradicts the assumption that $y$ is a vertex.  The claim follows.  In particular, $s_v$ is non-integral whenever the block corresponding to~$v$ is fractional.

Let $F\subseteq V(G)$ be the set of vertices whose blocks are fractional.  Define a graph~$T$ on vertex set~$F$ by putting $\{u,v\}\in E(T)$ if $\{u,v\}\in E(G)$ and the corresponding inequality is tight, namely
\begin{equation}\label{eq:tight}
 s_u+s_v=2.
\end{equation}
Note that a tight inequality cannot have exactly one endpoint in~$F$, since the sum of an integer and a non-integer cannot be~2.

We claim that $T$ is bipartite.  Otherwise, let $v_0v_1\cdots v_{2k}v_0$ be an odd cycle in~$T$.  The equations in \eqref{eq:tight} imply
\[
 s_{v_0}=s_{v_2}=\cdots=s_{v_{2k}}
 \qquad\text{and}\qquad
 s_{v_{2k}}=2-s_{v_0}.
\]
The final edge $v_{2k}v_0$ gives $2s_{v_0}=2$.  Hence $s_{v_0}=1$, contradicting the fact that $s_{v_0}$ is non-integral.

Now let $C$ be a connected component of~$T$, with bipartition $C=C^+\cupdot C^-$.  For every $v\in C$, perturb the unique fractional coordinate in its block by $+\delta$ if $v\in C^+$, and by $-\delta$ if $v\in C^-$.  The opposite perturbation is obtained by reversing the signs.  Every tight edge inequality incident with~$C$ has both endpoints in~$C$ and is preserved.  Every other affected edge inequality has positive slack.  Also, all the coordinates being changed lie strictly between 0 and~1.  Thus both perturbations are feasible for all sufficiently small $\delta>0$, again contradicting the assumption that $y$ is a vertex.  It follows that $F$ is empty, and hence every vertex of $Q(G)$ is binary.

It remains to show that the reduction mapping $G$ to $Q(G)$ is parsimonious.  Observe that there is a natural bijection $\phi:\{0,1\}^{V(G)\times\{1,2\}}\to V(\Hfour)^{V(G)}$ that interprets each binary vector $y\in\{0,1\}^{V(G)\times\{1,2\}}$ as a function $\sigma_y:V(G)\to V(\Hfour)$.  Explicitly, for each vertex $v\in V(G)$ we define $\sigma_y(v)=w$, where $w$ is the vertex of~$\Hfour$ encoded by the pair $(y_{v,1},y_{v,2})$.

In one direction, suppose that $y$ is a vertex of $Q(G)$.  We saw that $y$ is binary vector, so $\phi(y)$ is a function from $V(G)$ to $V(\Hfour)$.  It is easy to see that the construction of $Q(G)$ --- specifically the set of inequalities~\eqref{eq:edge} --- forces $\phi(y)$ to be a homomorphism.  This gives an injective map from vertices of $Q(G)$ to $\Hom(G,\Hfour)$.  In the other direction, any homomorphism  $\sigma\in\Hom(G,\Hfour)$ corresponds to a binary vector $y=\phi^{-1}(\sigma)$ and the construction of $Q(G)$ ensures that $y\in Q(G)$.  Any binary vector lying in a 0/1-polytope is a vertex of that polytope.  Thus the function $\phi^{-1}$ is an injective map from $\Hom(G,\Hfour)$ to vertices of $Q(G)$.
\end{proof}  

\end{document}
